\begin{helper}[Finite-prefix agreement]
If $x,y\in\mathsf{IncSeq}$ and the first $n$ values of $x$ form a proper
prefix of the range of $y$, then the first $n$ values of $x$ and $y$ agree
pointwise.
\end{helper}
\begin{proof}
Write $P_x=\{x(i)\mid i<n\}$.  By
\noderef{ProperPrefixSet}, choose a cutoff $y(j)$ such that
$P_x=\{y(i)\mid i<j\}$.  The order embeddings are injective, so both sides
have cardinalities $n$ and $j$, respectively.  Hence $n=j$.

We now prove $x(i)=y(i)$ for every $i<n$ by strong induction on $i$.  For
$i=0$, both $x(0)$ and $y(0)$ belong to the common finite set.  Every value
of either increasing enumeration is at least its zeroth value, so each is
at most the other; therefore they are equal.  For the induction step,
membership in the common set gives indices $k,j<n$ with
$x(k)=y(i+1)$ and $y(j)=x(i+1)$.  If $y(i+1)<x(i+1)$, strict increase of
$x$ gives $k<i+1$.  If $k=i$, the induction hypothesis gives
$x(i)=y(i)$ and hence $y(i)=y(i+1)$, contradicting strict increase; if
$k<i$, the induction hypothesis gives the same contradiction.  Thus
$x(i+1)\le y(i+1)$.  The symmetric argument gives the reverse inequality,
and antisymmetry gives equality.  This proves the stated prefix agreement.
\end{proof}
